\documentclass[10pt,a4paper]{amsart}
\usepackage{amsmath,amssymb,amsthm,graphicx}
% Retained minimal preamble for the Stacks Project chapter varieties.tex.
% Verbatim excerpt of the chapter preamble: Xy-pic setup (native preamble lines 14-18), the native theorem
% declarations with the mathematical macros (native preamble lines 154-217), and the project own no-op
% definitions of the reference/slogan/history display environments (native wrapper.tex lines 1-3) that the
% counted statement uses for its slogan. Copied mechanically; nothing added, removed or redefined.
% For commutative diagrams we use Xy-pic
\usepackage[all]{xy}

% We use 2cell for 2-commutative diagrams.
\xyoption{2cell}
% Theorem environments.
%
\theoremstyle{plain}
\newtheorem{theorem}[subsection]{Theorem}
\newtheorem{proposition}[subsection]{Proposition}
\newtheorem{lemma}[subsection]{Lemma}

\theoremstyle{definition}
\newtheorem{definition}[subsection]{Definition}
\newtheorem{example}[subsection]{Example}
\newtheorem{exercise}[subsection]{Exercise}
\newtheorem{situation}[subsection]{Situation}

\theoremstyle{remark}
\newtheorem{remark}[subsection]{Remark}
\newtheorem{remarks}[subsection]{Remarks}

\numberwithin{equation}{subsection}

% Macros
%
\def\lim{\mathop{\mathrm{lim}}\nolimits}
\def\colim{\mathop{\mathrm{colim}}\nolimits}
\def\Spec{\mathop{\mathrm{Spec}}}
\def\Hom{\mathop{\mathrm{Hom}}\nolimits}
\def\Ext{\mathop{\mathrm{Ext}}\nolimits}
\def\SheafHom{\mathop{\mathcal{H}\!\mathit{om}}\nolimits}
\def\SheafExt{\mathop{\mathcal{E}\!\mathit{xt}}\nolimits}
\def\Sch{\mathit{Sch}}
\def\Mor{\mathop{\mathrm{Mor}}\nolimits}
\def\Ob{\mathop{\mathrm{Ob}}\nolimits}
\def\Sh{\mathop{\mathit{Sh}}\nolimits}
\def\NL{\mathop{N\!L}\nolimits}
\def\CH{\mathop{\mathrm{CH}}\nolimits}
\def\proetale{{pro\text{-}\acute{e}tale}}
\def\etale{{\acute{e}tale}}
\def\QCoh{\mathit{QCoh}}
\def\Ker{\mathop{\mathrm{Ker}}}
\def\Im{\mathop{\mathrm{Im}}}
\def\Coker{\mathop{\mathrm{Coker}}}
\def\Coim{\mathop{\mathrm{Coim}}}

% Boxtimes
%
\DeclareMathSymbol{\boxtimes}{\mathbin}{AMSa}{"02}

%
% Macros for moduli stacks/spaces
%
\def\QCohstack{\mathcal{QC}\!\mathit{oh}}
\def\Cohstack{\mathcal{C}\!\mathit{oh}}
\def\Spacesstack{\mathcal{S}\!\mathit{paces}}
\def\Quotfunctor{\mathrm{Quot}}
\def\Hilbfunctor{\mathrm{Hilb}}
\def\Curvesstack{\mathcal{C}\!\mathit{urves}}
\def\Polarizedstack{\mathcal{P}\!\mathit{olarized}}
\def\Complexesstack{\mathcal{C}\!\mathit{omplexes}}
% \Pic is the operator that assigns to X its picard group, usage \Pic(X)
% \Picardstack_{X/B} denotes the Picard stack of X over B
% \Picardfunctor_{X/B} denotes the Picard functor of X over B
\def\Pic{\mathop{\mathrm{Pic}}\nolimits}
\def\Picardstack{\mathcal{P}\!\mathit{ic}}
\def\Picardfunctor{\mathrm{Pic}}
\def\Deformationcategory{\mathcal{D}\!\mathit{ef}}
\newenvironment{reference}{}{}
\newenvironment{slogan}{}{}
\newenvironment{history}{}{}
\title{Galois-invariant closed subsets of the base change descend to the ground scheme}
\author{The Stacks Project Authors}
\date{stacks-project:varieties.tex:local-snapshot-20261011; GNU FDL 1.2 or later}
\begin{document}
\maketitle
\noindent\textit{Editorial note.} All delivered mathematical text is the author's own, copied from the pinned Stacks Project chapter file varieties.tex; no statement or proof step was rewritten, reordered or removed.\par\medskip
\noindent\textit{Editorial note.} The counted statement and its proof are self-contained with respect to the pinned file: the proof cites no result of the project, so this task carries no external source references and no reference stubs.\par\medskip
\noindent\textit{Editorial note.} The preamble is the chapter's own retained minimal preamble (sources/stacks-varieties/retained-preamble.tex), a verbatim excerpt of the native Xy-pic setup, of the native theorem declarations and mathematical macros, and of the project's own no-op reference/slogan/history display environments. The native class line and the xr-hyper externaldocument block are omitted because the assembler supplies a plain amsart wrapper; no macro was invented or redefined.\par\medskip
\noindent\textit{Editorial note.} The Galois-theoretic facts used at proof lines 921-927 (finiteness of the stage carrying a given element, and the invariants of a finite Galois extension) appear among the standard prerequisites rather than as context spans, because the pinned file proves neither.\par\medskip
\noindent\textit{Allowed standard prerequisites (editorial).}\begin{itemize}
\item Basic Galois theory: invariants of a finite Galois extension and the orbit polynomial. Used to see that the coefficients of P(T) = prod\_\{sigma in G\}(T - sigma(x)) lie in A; the construction of the polynomial, the passage to a finite stage and the radical argument are all retained.
\item Radical ideals and the correspondence between closed subsets and radical ideals. The closed subset T-bar is encoded by its radical ideal I-bar inside the base change of A to k-bar; the proof shows that I-bar equals the radical of the ideal generated by I in that base change, and reads off T = V(I).
\end{itemize}\medskip
\begin{lemma}
\label{lemma-closed-fixed-by-Galois}
Let $k$ be a field. Let $\overline{k}/k$ be a (possibly infinite)
Galois extension. Let $X$ be a scheme over $k$. Let
$\overline{T} \subset X_{\overline{k}}$ have the following properties
\begin{enumerate}
\item $\overline{T}$ is a closed subset of $X_{\overline{k}}$,
\item for every $\sigma \in \text{Gal}(\overline{k}/k)$
we have $\sigma(\overline{T}) = \overline{T}$.
\end{enumerate}
Then there exists a closed subset $T \subset X$ whose inverse image
in $X_{\overline{k}}$ is $\overline{T}$.
\end{lemma}

\begin{proof}
This lemma immediately reduces to the case where $X = \Spec(A)$
is affine. In this case, let $\overline{I} \subset A \otimes_k \overline{k}$
be the radical ideal corresponding to $\overline{T}$.
Assumption (2) implies that $\sigma(\overline{I}) = \overline{I}$
for all $\sigma \in \text{Gal}(\overline{k}/k)$.
Pick $x \in \overline{I}$. There exists a finite Galois extension
$k'/k$ contained in $\overline{k}$ such that $x \in A \otimes_k k'$.
Set $G = \text{Gal}(k'/k)$. Set
$$
P(T) = \prod\nolimits_{\sigma \in G} (T - \sigma(x)) \in (A \otimes_k k')[T]
$$
It is clear that $P(T)$ is monic and is actually an element of
$(A \otimes_k k')^G[T] = A[T]$ (by basic Galois theory).
Moreover, if we write $P(T) = T^d + a_1T^{d - 1} + \ldots + a_d$
the we see that $a_i \in I := A \cap \overline{I}$.
Combining $P(x) = 0$ and $a_i \in I$ we find
$x^d = - a_1 x^{d - 1} - \ldots  - a_d \in I(A \otimes_k \overline{k})$.
Thus $x$ is contained in the radical of $I(A \otimes_k \overline{k})$.
Hence $\overline{I}$ is the radical of $I(A \otimes_k \overline{k})$ and
setting $T = V(I)$ is a solution.
\end{proof}

\end{document}
